\centerline{\bf \Huge Арифметика остатков}
\centerline{\bf \Huge }

\textbf{Определение.} Пусть $a$ и $b$ --- два целых числа, причем $b > 0$. Если $a = bq + r, где 0 \leq  r < b,$ то говорят, что $a$ дает при делении на $b$ неполное частное $q$ и остаток $r$.

Пример: $100 = 7 \cdot 14 + 2; (q = 14, r = 2) , \quad
- 50 = 7 \cdot (-8) + 6; (q = -8, r = 6).$

\begin{flushright}
        \scriptsize{}
\end{flushright}

\problem{1}
Число $x$ при делении на 12 дает остаток 7. Какой остаток получится при делении числа $x$ на 2; 3; 4; 6?

\problem{2}
Одно из целых чисел при делении на 7 дает остаток 5, а другое дает остаток 4. Какой остаток получится при делении на 7 их произведения?

\problem{3}
Докажите, что если числа $a$ и $b$ не делятся на $3$, но дают при делении на $3$ одинаковые остатки, то число $ab - 1$ делится на $3$.

\problem{4}
Докажите, что если \(a_1\) даёт остаток \(r_1\) при делении на \(m\), а \(a_2\) даёт остаток \(r_2\) при делении на \(m\), то
\begin{enumerate}
    \item[a)] число \(a_1 + a_2\) при делении на \(m\) даёт тот же остаток, что и число \(r_1 + r_2\);
    \item[б)] число \(a_1 - a_2\) при делении на \(m\) даёт тот же остаток, что и число \(r_1 - r_2\);
    \item[в)] число \(a_1 a_2\) при делении на \(m\) даёт тот же остаток, что и число \(r_1 r_2\);
    \item[г)] число \(a_1^n\) при делении на \(m\) даёт тот же остаток, что и число \(r_1^n\).
\end{enumerate}

\problem{5}
Найдите остаток от деления:

\begin{enumerate}
    \item[a)] $7^{123}$ на $9$
    \item[б)] $(-11)^{456}$ на $7$
    \item[в)] $(-121)^{2031}$ на $7$
\end{enumerate}

\problem{6}
Числа $p$, $p + 10$ и $p + 14$ --- простые. Найти все возможные значения числа $p$.